%%%%%%%%%%  scan idx 329  |  folio 316  %%%%%%%%%%
\(F^p\Lambda^n(E)=\Lambda^n(E)\), et si \(p<0\), \(F^p\Lambda^n(E)=0\). Ceci définit bien une filtration finie sur \(\Lambda^n(E)\).

On définit un homomorphisme~:
\[
  f_p:\Lambda^{n-p}(E')\otimes_A\Lambda^p(E'')
  \longrightarrow F^p\Lambda^n(E)/F^{p-1}\Lambda^n(E)
\]
en posant \(f_p(x'_1\wedge\cdots\wedge x'_{n-p}\otimes x''_1\wedge\cdots\wedge x''_p)=\operatorname{cl}_{\operatorname{Gr}^p}\bigl(u(x'_1)\wedge\cdots\wedge u(x'_{n-p})\wedge x_1\wedge\cdots\wedge x_p\bigr)\) qui ne dépend visiblement pas du choix des \(x_i\) relevant les sections données \(x''_i\) de \(E''\). Si la suite exacte est localement scindée, \(f_p\) est un isomorphisme, car localement \(\Lambda^n(E)\simeq\bigoplus_{p+q=n}\Lambda^p(E')\otimes_A\Lambda^q(E'')\).

Dans le cas ou \(E',E,E''\) sont localement facteurs directs de Modules libres de type fini, les \(\operatorname{Gr}^p\Lambda^n(E)\) sont localement facteurs directs de libres de type fini, donc les \(F^p\Lambda^n(E)\) également~; on peut donc écrire dans \(K\)~:
\(\lambda^n(E)=\sum_{p=0}^{n}\operatorname{cl}(\operatorname{Gr}^p\Lambda^n(E))=\sum_{p=0}^{n}\lambda^{n-p}(E')\cdot\lambda^p(E'')\), d'où l'additivité.

Soient \((T,A)\), \((T',A')\) deux topos commutativement annelés, et \(f:(T,A)\longrightarrow(T',A!)\) un morphisme de topos annelés. On sait (Exp. IV, 2.7) que \(f\) définit un morphisme d'anneau sur les \(K\) correspondants, soit \(f^*:K^{\bullet}_{\naif}(T',A')\longrightarrow K^{\bullet}_{\naif}(T,A)\). Comme l'image inverse commute aux puissances extérieures, le morphisme \(f^*\) correspondant sur les \(K\) est un \(\lambda\)-homomorphisme pour la structure qui vient d'être définie.

Signalons deux cas particuliers :
\begin{itemize}[label=-]
\item Soit \(X\) un schéma~; \(K^{\bullet}_{\naif}(X)\) est ainsi muni d'une structure de pré-\(\lambda\)-anneau; c'èst ce cas qui servira dans la suite du séminaire, et il peut être considéré comme « modèle » pour la théorie générale des 
